\section{Related Work}
\textbf{Classical Methods and Penalty Terms.} Early network pruning works can be 
broadly categorized into penalty-based and saliency-based approaches \cite{reed1993pruning}. 
Penalty methods augment the loss objective with sparsity-enforcing terms, such as 
L\_0 or L\_1 norms \cite{chauvin1989backpropagation, janowsky1989pruning, weigend1991generalization}. 
Saliency methods, on the other hand, score connection importance based on how much 
their removal affects model performance \cite{lecun1989optimal}. Classical saliency 
metrics include local weight or neuron sensitivity \cite{mozer1989using, karnin1990simple}, 
second-order Taylor expansions of the loss function, and diagonal or full Hessian 
matrices (e.g., Optimal Brain Damage \cite{lecun1989optimal} and Optimal Brain Surgeon \cite{hassibi1993second}) 
to prune and analytically update remaining weights without retraining. Although theoretically rigorous, calculating Hessian inverses is expensive for modern, large-scale architectures.

\textbf{Modern Post-Hoc & Dynamic Sparsity.} To bypass second-order derivative computation, magnitude-based pruning emerged as the standard post-hoc approach \cite{thimm1995evaluating, strom1997sparse}.
Han et al. \cite{han2015learning} popularized a three-step pipeline, pre-training, 
magnitude-based thresholding, and retraining. This was later combined with trained 
quantization and Huffman coding in Deep Compression \cite{han2016deep}. This post-hoc 
paradigm motivated the Lottery Ticket Hypothesis \cite{frankle2019lottery}, showing 
that dense, randomly initialized networks contain sparse subnetworks capable of 
learning in isolation when reset to their initial weights. To avoid train-prune-retrain cycles, static methods prune the network once at 
initialization using connection sensitivity (SNIP) \cite{lee2019snip}, gradient 
flow preservation (GraSP) \cite{wang2020grasp}, or data-agnostic synaptic flow 
conservation (SynFlow) \cite{tanaka2020pruning}. Alternatively, dynamic sparse 
training methods continuously optimize network parameters and connectivity budgets 
during training (SET \cite{mocanu2018scalable}; DeepR \cite{bellec2018deep}), or dynamically 
allocate layer-wise density using gradient and magnitude information (RigL) \cite{evci2020rigging}. 
However, these sparse training methods still rely on pre-specified, rigid layer-wise 
budgets, limiting the model's capacity to adapt organically.

\textbf{Structured Compression & Post-Training Optimization.} Structured pruning physically removes grouped parameters such as channels or filters. This yields direct hardware acceleration without specialized libraries \cite{sze2017efficient}. To handle complex inter-layer coupling, DepGraph \cite{fang2023depgraph} models layer dependencies as connected graphs to automate structural pruning. Post-training pruning methods compress trained models using small calibration datasets \cite{frantar2023sparsegpt}. WoodFisher \cite{singh2020woodfisher} uses empirical Fisher information to estimate Hessians, while methods like CBS \cite{yu2022combinatorial} and CHITA \cite{benbaki2023chita} use fast sparse regression to avoid Hessian inversion entirely.


\textbf{Large-Scale & Activity-Dependent Pruning.} Pruning billion-parameter models requires one-shot compression without retraining \cite{frantar2023sparsegpt}. Methods like SparseGPT \cite{frantar2023sparsegpt} and Wanda \cite{sun2024wanda} compress LLMs using weight reconstruction or weight-activation products. Closer to neurobiology, activity-dependent methods prune connections based on local signal perturbation \cite{dekhovich2023nnrelief, bingham2025fine}. In contrast to static metrics, post-hoc retraining, or fixed layer budgets, our proposed Dynamic Activity-Dependent Pruning (DADP) uses a reverse Hebbian mechanism to organically regulate layer capacity and induce structured neuron death during a single training pass.
